\chapter{Glosario de acrónimos}
\label{chap:acronimos}

%%%%%%%%%%%%%%%%%%%%%%%%%%%%%%%%%%%%%%%%%%%%%%%%%%%%%%%%%%%%%%%%%%%%%%%%%%%%%%%%

\begin{description}
  \item[API] \emph{Application Programming Interface}. Interfaz de programación
        de aplicaciones.
  \item[AST] \emph{Abstract Syntax Tree}. Árbol de sintaxis abstracta generado
        por el analizador sintáctico del intérprete Python a partir del código
        fuente.
  \item[CESGA] Centro de Supercomputación de Galicia.
  \item[CG] \emph{Conjugate Gradient}. Benchmark NAS que evalúa operaciones
        sobre matrices dispersas mediante el método del gradiente conjugado.
  \item[CPU] \emph{Central Processing Unit}. Unidad central de procesamiento.
  \item[DSL] \emph{Domain-Specific Language}. Lenguaje específico de dominio.
  \item[EP] \emph{Embarrassingly Parallel}. Benchmark NAS que evalúa la
        generación de pares gaussianos mediante el método de Box-Muller.
  \item[FFT] \emph{Fast Fourier Transform}. Transformada rápida de Fourier,
        algoritmo eficiente para calcular la transformada discreta de Fourier.
  \item[FT] \emph{Fourier Transform}. Benchmark NAS que evalúa la
        transformada discreta de Fourier tridimensional sobre un array
        complejo.
  \item[GIL] \emph{Global Interpreter Lock}. Mecanismo de bloqueo global del
        intérprete CPython que impide la ejecución simultánea de hilos Python
        sobre el bytecode del intérprete.
  \item[HPC] \emph{High Performance Computing}. Computación de altas
        prestaciones.
  \item[IS] \emph{Integer Sort}. Benchmark NAS que evalúa la ordenación de
        enteros mediante un algoritmo de tipo bucket sort.
  \item[MG] \emph{Multi-Grid}. Benchmark NAS que evalúa operadores de
        restricción, prolongación e interpolación sobre mallas en múltiples
        niveles.
  \item[Mop/s] Millones de operaciones por segundo. Métrica de rendimiento
        empleada por los NAS Parallel Benchmarks.
  \item[NAS] \emph{NASA Advanced Supercomputing}. División de la NASA que
        desarrolló los NAS Parallel Benchmarks.
  \item[NPB] \emph{NAS Parallel Benchmarks}. Conjunto de programas de
        referencia para la evaluación de sistemas de computación paralela.
  \item[OpenMP] \emph{Open Multi-Processing}. Estándar de programación
        paralela basado en directivas de compilador para sistemas de memoria
        compartida.
  \item[PEP] \emph{Python Enhancement Proposal}. Documento de propuesta de
        mejora del lenguaje Python.
\end{description}

%%%%%%%%%%%%%%%%%%%%%%%%%%%%%%%%%%%%%%%%%%%%%%%%%%%%%%%%%%%%%%%%%%%%%%%%%%%%%%%%
